% Problem 1 source: ne630_problems/lesson_02/statement.tex @ e1fba66, lines 89--153.
% Original problem and solution language retained.
\noindent{\bf Problem 1 [FNRP 1.9]}. A reactor operates at a power of $10^3$ MW(t) for 1 year.  Calculate the power from decay heat
\begin{enumerate}[label=(\alph*)]
\item 1 day following shutdown
\item 1 month following shutdown
\item 1 year following shutdown
\item Repeat a, b, and c, assuming only one month of operation, and compare results.
\end{enumerate}



\vspace{0.5cm}
\ifsolved
{
\color{purple}

\noindent{\bf SOLUTION}:

Equation (1.27) provides the Wigner-Way formula for decay heat:
\begin{equation*}
 P_d(t) = 0.0622 P_0 (t^{-0.2} - (t_0 + t)^{-0.2}) \, .
\end{equation*}
For all parts, $P_0 = 10^3$ \mega\watt. For (a)--(c), $t_0 = 1~\text{y} = 31536000\second$.
Then, for (a), $t = 24\times 3600 = 86400\second$ and
\begin{equation*}
 \begin{split}
   P_d(86400)
    &= 0.0622 \times 10^3 \times (86400^{-0.2} - (31536000+86400)^{-0.2}) \approx \boxed{4.44~\mega\watt} \\
 \end{split}
\end{equation*}
For (b), assume that an average month has 30 days so that
$t\approx 30\times24\times 3600=2592000$ and
\begin{equation*}
 \begin{split}
   P_d(2592000)
    &= 0.0622 \times 10^3 \times (2592000^{-0.2} - (31536000+2592000)^{-0.2}) \approx \boxed{1.31~\mega\watt} \\
 \end{split}
\end{equation*}
For (c), $t=t_0$ and
\begin{equation*}
 \begin{split}
   P_d(31536000)
    &= 0.0622 \times 10^3 \times (31536000^{-0.2} \\
    &- (31536000+31536000)^{-0.2}) \approx \boxed{0.255~\mega\watt} \\
 \end{split}
\end{equation*}
For (d), $t_0 = 30\times 24\times 3600 =2592000\second$ and
\begin{equation*}
 \begin{split}
   P_d(86400)
    &= 0.0622 \times 10^3 \times (86400^{-0.2} - (2592000+86400)^{-0.2}) \approx \boxed{3.182~\mega\watt} \\
   P_d(2592000)
    &= 0.0622 \times 10^3 \times (2592000^{-0.2} - (2592000+2592000)^{-0.2}) \approx \boxed{0.420\mega\watt} \\
   P_d(31536000)
    &= 0.0622 \times 10^3 \times (31536000^{-0.2} \\
    &- (2592000+31536000)^{-0.2}) \approx \boxed{0.031~\mega\watt} \\
 \end{split}
\end{equation*}

The primary impact of a shorter $t_0$ is that fewer radioactive fission products are generated and, thus, the decay heat is substantially reduced.

}
\newpage
\else
% NOTHING
\fi

% Problem 2 assignment source: ne630/homework/markdown/HW02.md.
% The Fall 2026 prompt substantially expands the older lesson_02 problem.
\noindent{\bf Problem 2 [Lewis Eq.~(1.24); adapted from DHNRA 2.19]}.  Lewis
states that more than 40 fission-fragment pairs can emerge from the fission of
\isoA{U}{235}.  Equation~(1.24) gives one example, \isoA{Xe}{140} and
\isoA{Sr}{94}.  Find one additional representative fission channel published
for thermal-neutron-induced fission of \isoA{U}{235}.  Your source must identify
both fragments and the prompt-neutron multiplicity; cite it.

\begin{enumerate}[label=(\alph*)]
\item Write the complete reaction, verify its nucleon-number and charge balance,
and compute its $Q$ value using tabulated atomic masses.
\item For either Lewis's fragment pair or the pair you selected, compare each
fragment's neutron-to-proton ratio with that of a stable isobar having the same
mass number.
\item Explain why fission fragments are generally neutron rich.
\item Predict the principal decay mode of each fragment and explain how that
decay moves the products toward stability.
\item Identify the particles emitted promptly during fission and those emitted
later as the fission products decay.  Which contribution to the fission energy
is effectively unrecoverable for power production?
\end{enumerate}

\begin{adaptedsolution}
% The introductory sentence and calculation structure below follow the legacy
% solution.  The published Sr-96/Xe-138 pair and expanded discussion are new.
Basic Googling can yield several examples.  In addition, there are computer
codes that can simulate the fission event using tabulated data.  Here, I use a
channel identified in the prompt gamma-coincidence measurements of
Mukhopadhyay et al.~\cite{mukhopadhyay2012}:
\[
 \isoAZ{U}{235}{92}+\isoAZ{n}{1}{0}
 \longrightarrow
 \isoAZ{Sr}{96}{38}+\isoAZ{Xe}{138}{54}
 +2\,\isoAZ{n}{1}{0}.
\]
The nucleon-number balance is
\[
 235+1=96+138+2=236,
\]
and the charge balance is
\[
 92+0=38+54+2(0)=92.
\]
Thus, the reaction conserves both nucleon number and charge.  Using AME2020
atomic masses~\cite{wang2021}, the corresponding $Q$ value is
\begin{align*}
 Q
 &=931.494\left[(235.043928117+1.00866491590)\right.\\
 &\hspace{2.7cm}\left.
 -(95.921719045+137.914146268+2(1.00866491590))\right]\ \mathrm{MeV}\\
 &\approx \boxed{185.74\ \mathrm{MeV}}.
\end{align*}

For the selected products,
\[
 \left(\frac{N}{Z}\right)_{\isoA{Sr}{96}}
 =\frac{58}{38}\approx1.526,
 \qquad
 \left(\frac{N}{Z}\right)_{\isoA{Xe}{138}}
 =\frac{84}{54}\approx1.556.
\]
For comparison, stable isobars \isoA{Mo}{96} and \isoA{Ba}{138} have
\[
 \left(\frac{N}{Z}\right)_{\isoA{Mo}{96}}
 =\frac{54}{42}\approx1.286,
 \qquad
 \left(\frac{N}{Z}\right)_{\isoA{Ba}{138}}
 =\frac{82}{56}\approx1.464.
\]
Both fission products therefore have larger neutron-to-proton ratios than the
corresponding stable isobars.

The compound nucleus \isoA{U}{236}$^*$ contains substantially more neutrons
than protons.  During scission, its nucleons are divided between two fragments
on a time scale too short for those fragments to move directly to the valley
of stability.  Prompt-neutron evaporation removes part of the excess and part
of the fragments' excitation energy, but the residual fragments generally
remain neutron rich.  They subsequently approach stability through beta decay.

The principal decay mode of both selected fragments is beta-minus decay:
\[
 n\longrightarrow p+e^-+\overline{\nu}_e.
\]
A beta-minus decay leaves $A$ unchanged, increases $Z$ by one, and decreases
$N$ by one.  It therefore lowers $N/Z$ and moves a neutron-rich isobar toward
stability.  Representative chains are
\[
 \isoA{Sr}{96}\xrightarrow{\beta^-}\isoA{Y}{96}
 \xrightarrow{\beta^-}\isoA{Zr}{96}
\]
and
\[
 \isoA{Xe}{138}\xrightarrow{\beta^-}\isoA{Cs}{138}
 \xrightarrow{\beta^-}\isoA{Ba}{138}\;(\text{stable}).
\]
Excited daughter nuclei can also emit decay gamma rays, and some sufficiently
neutron-rich fission products have beta-delayed-neutron branches.

Immediately after fission, the system contains two energetic heavy fragments,
prompt neutrons---two in the representative channel above---and prompt gamma
rays.  At later times, the fission products emit beta particles, electron
antineutrinos, decay gamma rays, and, for some precursors, delayed neutrons.
The two primary fragments carry roughly $80\%$ of the approximately
$200\ \mathrm{MeV}$ released in fission.  Fragment and beta-particle energy is
deposited over short ranges, while neutron and gamma-ray energy is ultimately
deposited elsewhere in the reactor and its shielding, apart from leakage.  The
antineutrinos interact so weakly that their energy is effectively unrecoverable
for power production.
\end{adaptedsolution}

\newpage

% Problem 3 assignment source: ne630/homework/markdown/HW02.md.
% No matching legacy solution was found; the solution is new for Fall 2026.
\noindent{\bf Problem 3 [Lewis Eqs.~(1.21)--(1.23)]}.  An idealized
prompt-neutron chain begins with $n_0$ neutrons.  Each neutron generation takes
an average time $\ell$, and the multiplication factor $k$ is constant.

\begin{enumerate}[label=(\alph*)]
\item Beginning with $n_{i+1}=kn_i$, derive
\[
 n(t)=n_0 k^{t/\ell}.
\]
\item For $k\approx1$, show that
\[
 n(t)\approx n_0\exp\left(\frac{k-1}{\ell}t\right).
\]
\item Classify systems having $k=0.995$, $k=1.000$, and $k=1.005$ as
subcritical, critical, or supercritical, and describe the long-term behavior
of the neutron population in each system.
\item A reactor initially produces $1.25~\mathrm{MW(t)}$, with $k=1.010$ and
$\ell=10^{-5}~\mathrm{s}$.  Assuming $P(t)\propto n(t)$, estimate the time
required to reach $2.00~\mathrm{MW(t)}$ using both the exact and approximate
expressions.
\end{enumerate}

\begin{draftsolution}
After $i$ neutron generations,
\[
 n_i=n_0k^i.
\]
If each generation takes an average time $\ell$, then $i=t/\ell$.
Consequently,
\[
 \boxed{n(t)=n_0k^{t/\ell}}.
\]

Writing the power as an exponential gives
\[
 n(t)=n_0\exp\left(\frac{\ln k}{\ell}t\right).
\]
For $k$ near one,
\[
 \ln k=\ln\!\left[1+(k-1)\right]
       =(k-1)-\frac{(k-1)^2}{2}+\cdots\approx k-1.
\]
Therefore,
\[
 \boxed{n(t)\approx n_0\exp\left(\frac{k-1}{\ell}t\right)}.
\]

Within this idealized source-free model,
\[
\begin{array}{ccl}
 k=0.995<1&\Longrightarrow&\text{subcritical; }n(t)\text{ decreases toward zero},\\[1mm]
 k=1.000&\Longrightarrow&\text{critical; }n(t)=n_0\text{ remains constant},\\[1mm]
 k=1.005>1&\Longrightarrow&\text{supercritical; }n(t)\text{ increases exponentially}.
\end{array}
\]
The unlimited growth in the last line is a consequence of the simplified
model; real systems eventually experience feedback and other effects not
included here.

Because $P(t)\propto n(t)$,
\[
 \frac{P(t)}{P(0)}=\frac{2.00}{1.25}=1.60.
\]
Using the exact generation expression,
\begin{align*}
 t_{\mathrm{exact}}
 &=\ell\frac{\ln(1.60)}{\ln k}\\
 &=(10^{-5}\ \mathrm{s})\frac{\ln(1.60)}{\ln(1.010)}
 \approx\boxed{4.7235\times10^{-4}\ \mathrm{s}}
 =\boxed{0.47235\ \mathrm{ms}}.
\end{align*}
Using the $k\approx1$ approximation,
\begin{align*}
 t_{\mathrm{approx}}
 &=\ell\frac{\ln(1.60)}{k-1}\\
 &=(10^{-5}\ \mathrm{s})\frac{\ln(1.60)}{0.010}
 \approx\boxed{4.7000\times10^{-4}\ \mathrm{s}}
 =\boxed{0.47000\ \mathrm{ms}}.
\end{align*}
The approximation is lower by about $2.35\times10^{-6}\ \mathrm{s}$, or
$0.497\%$, because $\ln(1.010)<0.010$.
\end{draftsolution}

\clearpage
